\section{Transport Layer}
\subsection{Connection Control}
\subsubsection*{Grundidee}
In Netzwerken gibt es \fatsf{zwei Arten von Diensten:}
\begin{table}[H]
    \centering
    \rowcolors{2}{\thepagecolor}{fgcolor!10!\thepagecolor}
    \renewcommand{\arraystretch}{1.2}
    \begin{tabular}{lll}
        \toprule
        \fatsf{Art} & \fatsf{Erklärung} & \fatsf{Besipiel (Internet)}\\
        \midrule
        \fatsf{Connection-oriented} (COTS) & Verbindungsaufbau vor dem Datentransfer & \fatsf{TCP} (z.B. bei Web, E-Mail)\\
        \fatsf{Connectionless} (CLNS) & Direkter Versand ohne Verbindung & \fatsf{UDP} (z.B. bei VoIP, DNS)\\
        \bottomrule 
    \end{tabular}
\end{table}
Im folgenden Fokus: \fatsf{Connection-oriented Transport Services (COTS)}, wie sie z.B. bei TCP verwendet werden
\subsubsection*{Drei Phasen einer Verbindung (OSI-Terminologie)}
\begin{table}[H]
    \centering
    \rowcolors{2}{\thepagecolor}{fgcolor!10!\thepagecolor}
    \renewcommand{\arraystretch}{1.2}
    \begin{tabular}{lp{5cm}p{8cm}}
        \toprule
        \fatsf{Phase} & \fatsf{Beschreibung} & \fatsf{OSI Primitive (Beispiel)}\\
        \midrule
        \fatsf{1.Connect (Aufbau)} & Eine Verbindung wird hergestellt & T-Connect.Request, T-Connect.Indication, T-Connect.Response, T-Connect.Confirmation\\
        \fatsf{2. Data (Transofer)} & Daten werden zwischen Endsystemen übertragen & T-Data.Request, T-Data.Indication\\
        \fatsf{3. Disconnect (Abbau)} & Die Verbindung wird ordentlich beendet & T-Disconnect.Request, T-Disconnect.Indication\\
        \bottomrule 
    \end{tabular}
\end{table}
\subsubsection*{Transport-Verbindung: Beteiligte}
\begin{figure}[h]
    \centering
    \includegraphics[width=0.8\linewidth]{Chapter_4_Images/Transport ConnectiSystem Connectivity.jpeg}
\end{figure}
Ein \fatsf{SAP (Service Access Point)} ist ein \fatsf{logischer Zugangspunkt}, über den eine Schicht \fatsf{einen Dienst} von der \fatsf{darunterliegenden Schicht} anfordert
\subsubsection*{Service-Primitives (OSI-Sicht)}
Diese sind \fatsf{abstrakte Befehle}, die zwischen Schichten ausgetauscht werden, um Dienste zu nutzen\\
\begin{figure}[h]
    \centering
    \includegraphics[width=0.5\linewidth]{Chapter_4_Images/Example Run.jpeg}
\end{figure}
\fatsf{Phase 1: Verbindung aufbauen (T-Connect)}
\begin{enumerate}
    \item \fatsf{T-Connect.Request}
    \begin{itemize}
        \item Wer? $\rightarrow$ Client (aktive Seite)
        \item Richtung $\rightarrow$ vom Transport Service User an die Transport-Schicht
        \item Bedeutung:
        \begin{itemize}
            \item Der Client sagt: \enquote{Ich möchte eine Verbindung zu einem bestimmten Ziel aufbauen}
        \end{itemize}
        \item \fatsf{Parameter:}
        \begin{itemize}
            \item \texttt{Destination Address:} z.B. IP-Adresse + Port des Servers (Zielanwendung)
            \item \texttt{Source Address:} z.B. IP-Adresse + Port des Clients (eigene Anwendung)
        \end{itemize}
    \end{itemize}
    \item \fatsf{T-Connect.Indication}
    \begin{itemize}
        \item Wer? $\rightarrow$ Server (passive Seite), ausgelöst durch Transport-Schicht
        \item Richtung $\rightarrow$ von Transport-Schicht an die Server-Anwendung
        \item Die Transport-Schicht informiert die Server-Anwendung: \enquote{Hey, da möchte jemand mit dir eine Verbindung aufbauen!}
        \item \fatsf{Wichtig:}
        \begin{itemize}
            \item \fatsf{Asynchron} - der Server \enquote{lauscht} mit einem listen()-Mechanismus
            \item Das passiert nicht durch aktiven Aufruf des Servers, sondern als \fatsf{Indication (Meldung)}
        \end{itemize}
    \end{itemize}
    \item \fatsf{T-Connect.Response}
    \begin{itemize}
        \item Wer? $\rightarrow$ Server (passive Seite)
        \item Richtung $\rightarrow$ von Server-Anwendung an Transport-Schicht
        \item Bedeutung:
        \begin{itemize}
            \item Der Server sagt: \enquote{Ja, ich akzeptiere die Verbindung}
        \end{itemize}
        \item Parameter:
        \begin{itemize}
            \item \texttt{Responding Address:} Adresse des Servers (kann benutzt werden, um Verbindung zu bestätigen und zu identifizieren)
        \end{itemize}
    \end{itemize}
    \item \fatsf{T-Connect.Confirmation}
    \begin{itemize}
        \item Wer? $\rightarrow$ Client (aktive Seite), ausgelöst durch Transport-Schicht
        \item Richtung $\rightarrow$ von Transport-Schicht an die Client-Anwendung
        \item Bedeutung:
        \begin{itemize}
            \item Die Transport-Schicht bestätigt dem Client: \enquote{Verbindung ist aufgebaut}
        \end{itemize}
        \item Optionaler Parameter:
        \begin{itemize}
            \item Responding Address: Adresse des Servers, um finale Details zu bestätigen
        \end{itemize}
    \end{itemize}
\end{enumerate}
\subsubsection*{Datenübertragung - T-Data}
\begin{itemize}
    \item \fatsf{Merkmale:}
    \begin{itemize}
        \item \fatsf{Unconfirmed service:} Es wird \fatsf{kein explizites Bestätigungs-Primitive (ACK)} auf der Service-Ebene gesendet
        \item Trotzdem: \fatsf{Das Protokoll (z.B. TCP)} sendet \fatsf{ACKs intern!}
    \end{itemize}
    \item \fatsf{Primitives:}
    \begin{itemize}
        \item T-Data.Request (userdata) - vom Sender aufgerufen
        \item T-Data.Indication (userdata) - beim Empfänger ausgelöst
    \end{itemize}
\end{itemize}
\clearpage
\subsubsection*{Verbindungsabbau - T-Disconnect}
Wieder ein 2-Wege-Prozess:
\begin{grayInfoBox}
    Client: T-Disconnect.Request\\
    Server: T-Disconnect.Indication
\end{grayInfoBox}
Das entspricht dem TCP \texttt{FIN}/\texttt{ACK}-Mechanismus, also:
\begin{itemize}
    \item Ein Ende sendet \texttt{FIN}
    \item Das andere bestätigt mit \texttt{ACK}
\end{itemize}
\begin{figure}[h]
    \centering
    \includegraphics[width=0.8\linewidth]{Chapter_4_Images/Connection Release.jpeg}
\end{figure}
\subsubsection*{Zustandsdiagramm für einen Transport Service Access Point}
\begin{figure}[h]
    \centering
    \includegraphics[width=0.8\linewidth]{Chapter_4_Images/State Diagram.jpeg}
\end{figure}
\subsubsection*{Beispiel: TCP-Fall durchspielen (vereinfacht)}
\begin{table}[H]
    \centering
    \rowcolors{2}{\thepagecolor}{fgcolor!10!\thepagecolor}
    \renewcommand{\arraystretch}{1.2}
    \begin{tabular}{llll}
        \toprule
        \fatsf{Phrase} & \fatsf{Zustand} & \fatsf{Service Primitive} & \fatsf{Bemerkung}\\
        \midrule
        Client startet & $1 \rightarrow 2$ & \texttt{T-Connect.req} & TCP sendet SYN\\
        Server erhält Anfrage & $1 \rightarrow 3$ & \texttt{T-Connect.ind} & TCP empfängt SYN\\
        Server antwortet & $3 \rightarrow 4$ & \texttt{T-Connect.rsp} & TCP sendet SYN+ACK\\
        Client bestätigt & $2 \rightarrow 4$ & \texttt{T-Connect.cnf} & TCP empfängt SYN+ACK + sendet ACK\\
        Beide senden Daten & 4 & \texttt{T-Data.req/ind} & TCP-Transfer läuft\\
        Verbindung wird beendet & $4 \rightarrow 1$ & \texttt{T-Disconnect.req/ind} & TCP sendet FIN, bekommt FIN\\
        \bottomrule 
    \end{tabular}
\end{table}
\subsubsection*{Errors during Connection Establishment}
Beim Aufbau einer Verbindung (Phrase: Connection Establishment) werden zwischen zwei Endpunkten \fatsf{Kontrollnachrichten} (TPDUs - Transport Protocol Data Units) ausgetauscht. Die wichtigsten:
\begin{itemize}
    \item \fatsf{CR} = Connection Request $\rightarrow$ vom Client gesendet
    \item \fatsf{CC} = Connection Confirm $\rightarrow$ vom Server gesendet
\end{itemize}
Wenn dabei Fehler passieren (z.B. Nachrichten gehen verloren oder werden doppelt empfangen), kann es zu \fatsf{Problemen im Verbindungsaufbau} kommen
\begin{figure}[h]
    \centering
    \includegraphics[width=0.8\linewidth]{Chapter_4_Images/Errors during Connection Establishment.jpeg}
\end{figure}
\subsubsection*{Fehlersituationen im Detail}
\begin{enumerate}
    \item \fatsf{Verlust von CR oder CC (Nachricht kommt nie an)}
    \begin{itemize}
        \item \texttt{T-Connect.req} $\rightarrow$ CR gesendet
        \item \fatsf{CR kommt beim Server nicht an $\rightarrow$ Server reagiert nicht}
        \item Lösung: \fatsf{Client startet nach Timeout neu}
    \end{itemize}
    \item \fatsf{Deppelte TPDUs (z.B. durch Netzwerkprobleme)}
    \begin{itemize}
        \item CR oder CC wird \fatsf{mehrmals empfangen}, z.B. durch Wiederholung nach Timeout
    \end{itemize}
    \fatsf{Beispiel:} 
    \begin{itemize}
        \item CR wird vom Client gesendet
        \item Server empfängt CR $\rightarrow$ sendet CC
        \item CC kommt beim Client nicht an $\rightarrow$ Client \fatsf{wiederholt} CR
        \item Server empfängt CR \fatsf{ein zweites Mal $\rightarrow$ nochmals} \texttt{T-connect.ind} \fatsf{an Anwendung}
    \end{itemize}
    \fatsf{Problem:} Anwendung bekommt zweimal die gleiche Verbindungsanfrage! Das darf \fatsf{nicht passieren}\\
    Die Anwendungsschicht \fatsf{soll vor diesen Fehlern geschützt} werden. Solche Duplikate sollen vom \fatsf{Transportprotokoll abgefangen} werden, \fatsf{bevor} sie zu Anwendung gelangen
\end{enumerate}
\subsubsection*{Three-Way Handshake}
\begin{itemize}
    \item \fatsf{Zweck:}
    \begin{itemize}
        \item Sicherer Verbindungsaufbau im Transport Layer (z.B. TCP)
        \item Verhindert, dass Verbindungen inkonsistent oder fehlerhaft aufgebaut werden
    \end{itemize}
    \item \fatsf{Ablauf:}
    \begin{enumerate}
        \item \fatsf{CR} (Connection Request) vom Client $\rightarrow$ Server
        \item \fatsf{CC} (Connection Confirm) vom Server $\rightarrow$ Client
        \item \fatsf{ACK} oder erste \fatsf{DT} (Data Transfer) vom Client $\rightarrow$ Server
    \end{enumerate}
    \fatsf{Verbindung ist erst nach allen 3 Schritten offiziell etabliert}
\end{itemize}
\begin{figure}[h]
    \centering
    \includegraphics[width=0.6\linewidth]{Chapter_4_Images/Three-Way Handshake.jpeg}
\end{figure}
\subsubsection*{Problem: Alte oder doppelte Pakete}
\begin{itemize}
    \item Netzwerk kann alte Verbindungsanfragen (CR/CC) später erneut zustellen (sog. Replay-Angriffe)
    \item Gefahr: eine \enquote{vergangene} Verbindung wird fälschich als neu interpretiert
\end{itemize}
\subsubsection*{Lösung: Sequence Numbers}
\begin{itemize}
    \item Jeder Teilnehmer sendet \fatsf{eine zufällige Nummer (SeqNum)} in CR/CC
    \item Die Gegenseite muss diese Nummer korrekt zurückspiegeln
    \item Verbindung wird nur akzeptiert, wenn die SeqNums \fatsf{zeitlich gültig und korrekt} sind
\end{itemize}
\fatsf{Verhindert den Aufbau auf Basis alter, duplizierter Nachrichten}
\begin{figure}[h]
    \centering
    \includegraphics[width=0.3\linewidth]{Chapter_4_Images/Sequence numbers.jpeg}
\end{figure}
\subsubsection*{Zusätzlicher Schutz}
\begin{itemize}
    \item \fatsf{Sequence Numbers dürfen sich nicht wiederholen}, solange alte Pakete noch im Netz sein könnten
    \begin{itemize}
        \item $\rightarrow$ Zeit > \enquote{Packet Lifetime} (oft durch \texttt{TTL} im IP-Layer begrenzt)
    \end{itemize}
    \item Bei sehr schnellen Verbindungen (z.B. Gbps):\\
    $\rightarrow$ \fatsf{Timestamps} zusätzlich verwenden
\end{itemize}
\begin{figure}[h]
    \centering
    \includegraphics[width=1\linewidth]{Chapter_4_Images/Old duplicate.jpeg}
\end{figure}
\clearpage
\subsubsection*{Connection Rejection (Verbindungsablehnung)}
Wenn ein Transportdienst \fatsf{eine eingehende Verbindungsanfrage (CR) nicht akzeptieren} will, wird sie \fatsf{aktiv zurückgewiesen}
\begin{itemize}
    \item \fatsf{Wie funktioniert das?}
    \begin{itemize}
        \item \fatsf{Client sendet:} \texttt{T-Connect.req} (Verbindungsaufbau starten)
        \item \fatsf{Server lehnt ab:} mit \texttt{T-Disconnect.ind} und einem \fatsf{DR} (Disconnect Request TPDU)
        \item Danach folgt evtl. ein \fatsf{DC} (Disconnect Confirm TPDU) zur Bestätigung
    \end{itemize}
    \item \fatsf{Mögliche Gründe für Ablehnung:}
    \begin{itemize}
        \item Transportdienstnutzer will Verbindung nicht (z.B. Anwendung sagt \enquote{Nein})
        \item Dienstanforderungen (z.B. QoS) können nicht erfüllt werden
    \end{itemize}
\end{itemize}
\begin{figure}[h]
    \centering
    \includegraphics[width=0.5\linewidth]{Chapter_4_Images/Connection Rejection.jpeg}
\end{figure}
\subsubsection*{Connection Release (Verbindungsabbau)}
\begin{itemize}
    \item \fatsf{Ziel:}
    \begin{itemize}
        \item Verbindung \fatsf{kontrolliert beenden}, nachdem \fatsf{alle Daten sicher übertragen werden}
    \end{itemize}
    \item \fatsf{Empfohlener Ablauf:}
    \begin{enumerate}
        \item Warten, bis alle gesendeten Daten bestätigt (ACKed) wurden
        \item Dann erst: Verbindung schließen
    \end{enumerate}
    \item \fatsf{Varianten:}
    \begin{itemize}
        \item \fatsf{Explizit:} durch \fatsf{DR/DC} (Disconnect Request + Confirm TPDU)
        \item \fatsf{Implizit:} Einfach durch Schließen der Verbindung im darunterliegenden Netzwerk
    \end{itemize}
\end{itemize}
\begin{figure}[h]
    \centering
    \includegraphics[width=0.8\linewidth]{Chapter_4_Images/Connection Release 2.jpeg}
\end{figure}
\subsubsection*{Problem beim Verbindungsabbau}
Wie kannst du dir sicher sein, dass dein Kommunikationsparter \fatsf{weiß}, dass \fatsf{du weißt}, dass \fatsf{er weiß}, dass alles angekommen ist?\\
$\rightarrow$ Genau das ist das \fatsf{Two-Army-Problem!}
\subsubsection*{Two-Army-Problem-Analogie}
\begin{itemize}
    \item Zwei Armeen wollen gleichzeitig angreifen, sonst verlieren sie
    \item Sie schicken Boten zueinander - aber Boten können \fatsf{verloren} gehen
    \item Sie wissen nie, ob die Nachricht \fatsf{wirklich angekommen} ist
    \item Ergebnis: \fatsf{keine sichere Koordination möglich}, wenn Nachrichten verloren gehen können
\end{itemize}
\begin{itemize}
    \item \fatsf{Gleiches Problem beim Connection Release!}
    \item Transport Layer kann \fatsf{nicht sicher wissen}, ob wirklich beide Seiten bereit sind zu schließen
\end{itemize}
\begin{figure}[h]
    \centering
    \includegraphics[width=0.8\linewidth]{Chapter_4_Images/Two Army Problem.jpeg}
\end{figure}
\subsubsection*{Praktische Lösung: 3-Way Handshake für Verbindungsabbau}
\begin{enumerate}
    \item Peer A sendet \fatsf{DR}
    \item Peer B antwortet auch mit \fatsf{DR} (ggf. inkl. ACK)
    \item A bestätigt, Verbindung kann geschlossen werden
\end{enumerate}
$\rightarrow$ Ist zwar \fatsf{100\% sicher}, aber:
\begin{itemize}
    \item In der Praxis funktioniert das \fatsf{gut genug}
    \item \fatsf{Timer} und \fatsf{Wiederholungen} helfen bie verlorenen DR/ACK
\end{itemize}
\subsubsection*{Zusätzliche Problemfälle:}
\begin{itemize}
    \item \fatsf{ACK verloren:} Timer läuft aus $\rightarrow$ erneutes DR senden
    \item \fatsf{2. DR verloren:} ebenfalls durch Timer abgefangen
    \item Wenn \fatsf{beide nachrichten verloren:} Verbindung wird trotzdem nach Timeout abgebrochen
\end{itemize}
\begin{figure}[h]
    \centering
    \includegraphics[width=1\linewidth]{Chapter_4_Images/Problem Cases for Connection Release.jpeg}
\end{figure}
\subsection{Flow Control}
\fatsf{Flow Control} bedeutet, dass ein \fatsf{Sender daran gehindert wird}, zu viele Daten zu schnell an einen \fatsf{langsamen Empfänger} zu senden. Ziel: Den \fatsf{Empfänger vor Überlastung} zu schützen!
\subsubsection*{Hintergrundproblem: Zuverlässige Übertragung ist schwer}
Auch wenn Pakete zuverlässig zugestellt werden sollen, können viele Dinge schieflaufen:
\begin{itemize}
    \item Pakete gehen verloren (z.B. durch Überlastung oder Bitfehler)
    \item Pakete kommen fehlerhaft an
    \item Der Empfänger verarbeitet Daten langsamer, als sie ankommen
\end{itemize}
\subsubsection*{Lösungsidee: ACKs und Retransmission}
Damit der Sender weiß, ob Daten korrekt ankamen:
\begin{itemize}
    \item Der Empfänger schickt ein \fatsf{ACK} (Bestätigung) für jedes empfangene Paket
    \item Kommt kein ACK $\rightarrow$ Sender \fatsf{sendet Paket erneut} (Der Sender muss eine Kopie des Pakets behalten, bis er das ACK sicher erhalten hat - nur dann kann er es löschen)
\end{itemize}
\subsubsection*{Aber: Wie lange soll der Sender ein Paket speichern?}
\begin{itemize}
    \item Solange, bis er ein ACK erhält
    \item Das braucht Speicher (= \fatsf{Buffer}) $\rightarrow$ Also \fatsf{nicht unbegrenzt machbar}
\end{itemize}
\subsubsection*{Flow Control - Warum wichtig?}
\begin{itemize}
    \item \fatsf{Problem:} Ein \fatsf{schneller Sender} kann einen \fatsf{langsamen Empfänger} überfordern, wenn er \fatsf{ungebremst Daten schickt}
    \item \fatsf{Ziel:} Empfänger muss \fatsf{Sender bremsen können}, wenn sein \fatsf{Eingangspuffer (Buffer)} voll ist
\end{itemize} 
\subsubsection*{Wo wird Flow Control gemacht?}
\begin{itemize}
    \item \fatsf{Link Layer:} begrenzt Datenrate zwischen zwei direkt verbundenen Geräten
    \item \fatsf{Transport Layer:} steuert den Datenfluss \fatsf{zwischen Anwendungen} auf verschiedenen Rechnern
    \begin{itemize}
        \item Komplizierter, da:
        \begin{itemize}
            \item TL-Pakete unterschiedlich groß sind
            \item Anwendungen (z.B. Webserver) sehr unterschiedlich reagieren
            \item Mehrere parallele Verbindungen gleichzeitig laufen können
        \end{itemize}
    \end{itemize}
\end{itemize}
\subsubsection*{Buffer Management}
Sender und Empfänger müssen mit \fatsf{Pufferplätzen} umgehen:
\begin{itemize}
    \item Der \fatsf{Empfänger} braucht Platz, um Pakete zwischenzuspeichern, bevor er sie verarbeitet
    \item Der \fatsf{Sender} muss wissen, \fatsf{wie viel Platz beim Empfänger verfügbar ist}
\end{itemize}
\fatsf{Zwei Optionen:}
\begin{enumerate}
    \item \fatsf{Optimistische Annahme:} Empfänger verarbeitet Pakete sofort $\rightarrow$ unrealistisch
    \item \fatsf{Guthabensystem (Credit-Based):}
    \begin{itemize}
        \item Empfänger sagt: \enquote{Ich habe noch \fatsf{X feie Buffer}}
        \item Sender darf \fatsf{nur so viele neue Pakete schicken}, wie \fatsf{Credits} vorhanden sind
    \end{itemize}
    $\rightarrow$ So kann der Empfänger den \fatsf{Datenfluss kontrollieren}, ohne überlastet zu werden
\end{enumerate}
\subsubsection*{Wie kommuniziert der Empfänger seine Kapazität?}
\begin{itemize}
    \item \fatsf{Explizit:} \enquote{Ich habe gerade 4 freie Plätze}
    \item \fatsf{Implizit:} Durch Ausbleiben von ACKs oder bestimmte TCP-Felder
\end{itemize}
\clearpage
\subsubsection{Alternating Bit Protocol (ABP)}
Das \fatsf{ABP} ist ein einfaches Protokoll zur \fatsf{sicheren (zuverlässigen)} und \fatsf{geordneten Datenübertragung} über ein unsicheres Netzwerk - also ein \fatsf{Sotp-and-Wait-Protokoll mit Fehlererkennung}
\subsubsection*{Grundidee}
\begin{enumerate}
    \item Der \fatsf{Sender} sendet \fatsf{ein Paket mit einer Bit-Nummer} (entweder 0 oder 1)
    \item Der \fatsf{Empfänger} erhält das Paket und schickt ein \fatsf{ACK mit der gleichen Bit-Nummer} zurück
    \item Der Sender wartet auf das ACK, \fatsf{bevor} er das nächste Paket sendet - das nächste mit \fatsf{wechselndem Bit} (0 $\leftrightarrow$ 1)
\end{enumerate}
Wenn kein ACK rechtzeitig kommt (Timeout), \fatsf{wird dasselbe Paket nochmal gesendet}
\subsubsection*{Warum nur 1 Bit?}
Weil immer \fatsf{nur ein Paket gleichzeitig} unterwegs ist. Das Bit dient zur Erkennung von \fatsf{Doppelübertragungen}
\begin{table}[H]
    \centering
    \rowcolors{2}{\thepagecolor}{fgcolor!10!\thepagecolor}
    \renewcommand{\arraystretch}{1.2}
    \begin{tabular}{p{7cm}p{8cm}}
        \toprule
        \fatsf{Situation} & \fatsf{Erklärung}\\
        \midrule
        Paket kommt \fatsf{neu} $\rightarrow$ Bitnummer \fatsf{neu} & $\rightarrow$ Empfänger verarbeitet \& bestätigt\\
        Paket kommt \fatsf{doppelt} $\rightarrow$ Bitnummer wie beim letzten Mal & $\rightarrow$ Empfänger \fatsf{verarbeitet es nicht nochmal} und schickt \fatsf{ACK (Bit = Bitnummer)}\\
        \bottomrule 
    \end{tabular}
\end{table}
\subsubsection*{Ist ABP effizient?}
Leider \fatsf{nein} - und das ist ein zentrales Lernziel:
\begin{figure}[h]
    \centering
    \includegraphics[width=0.7\linewidth]{Chapter_4_Images/Rechenbeispiel.jpeg}
\end{figure}
\begin{itemize}
    \item \fatsf{Entfernung:} New York $\leftrightarrow$ Frankfuhrt = 6.200 km
    \item \fatsf{Lichgeschwindigkeit in Glasfaser:} 300.000 km/s
    \item \fatsf{Propagation through the medium:}
    \[
        T_{prop} = \frac{6200km}{300000\frac{km}{s}} = 20ms
    \]
    \item \fatsf{Transmission Time:} 
    \[
        T_{trans} = \frac{8KBit}{10^{9}\frac{Bit}{s}} = 8 \mu s \text{ (assume 1 $\mu s$ for ACK)}
    \]
    \item \fatsf{Delay until ACK of packet received:} 40009 ms
    \item \fatsf{Utilization of medium:} 
    \[
        U_{Sender} = \frac{0,008}{40,009} = 0.0002
    \]
    \item And effectively around 25 kByte/s (in a 1 GBit/s medium)
\end{itemize}
Also:
\begin{itemize}
    \item Extrem ineffizient auf langen Strecken und bei hoher Bandbreite
    \item Aber einfach und sicher für kurze Distanzen (z.B. Embedded-Bereich)
\end{itemize}
\begin{table}[H]
    \centering
    \rowcolors{2}{\thepagecolor}{fgcolor!10!\thepagecolor}
    \renewcommand{\arraystretch}{1.2}
    \begin{tabular}{lp{4cm}p{4cm}p{4cm}}
        \toprule
        \fatsf{Konzept} & \fatsf{Was passiert?} & \fatsf{Vorteil} & \fatsf{Nachteil}\\
        \midrule
        Stop/Continue Flow Control & Empfänger schickt Stop/Continue Nachrichten & Einfach zu verstehen & Verzögerung kann Overflow verursachen\\
        Rate-based Flow Control & Empfänger gibt max. Datenrate an & Geringer Overhead, \enquote{smooth} & Kein Empfangs-ACK, Fehler schwierig\\
        Credit-Based Flow Control & Empfänger gibt explizites Guthaben (Pakete senden erlaubt) & Effizient, kombiniert mit ACKs & Komplexer, benötigt Fehlerkontrolle\\
        Sliding Window & Sender darf mehrere Pakete im Fenster senden & Sehr effizient, verbreitet & Verwaltung von Fenstergröße nötig\\
        \bottomrule 
    \end{tabular}
\end{table}
\clearpage
\subsection{Error Control}
\fatsf{Problem:} Nicht alle Übertragungsfehler werden erkannt
\begin{itemize}
    \item Bei Datenübertragungen können \fatsf{Bitfehler} auftreten (z.B. durch elektromagnetische Störungen)
    \item Man verwendet \fatsf{Prüfsummen (Checksums)}, um solche Fehler zu erkennen
    \item Aber: \fatsf{Es gibt immer ein Restrisiko}, dass ein Fehler \fatsf{nicht entdeckt wird}, wenn die falschen Daten zufälltig dieselbe Prüfsumme haben $\Rightarrow$ \fatsf{residual error}
\end{itemize}
\fatsf{Verhalten auf Empfängerseite bei erkannten Fehlern:}
\begin{itemize}
    \item Wenn ein Fehler erkannt wird (z.B. durch Prüfsumme in Layer 2 oder Layer 4)
    \begin{itemize}
        \item \fatsf{In den meisten Protokollen reagiert der Empfänger NICHT aktiv}
        \item Das fehlerhafte Paket wird \fatsf{einfach verworfen (drop)}
    \end{itemize}
    \item \fatsf{Warum wird einfach verworfen?}
    \begin{itemize}
        \item Gefahr: \fatsf{Ein fehlerfhaftes Headerfeld} (z.B. Zieladresse) könnte das Paket \fatsf{in eine falsche Richtung} weiterleiten $\Rightarrow$ \fatsf{besser verwerfen}
        \item Es gibt \fatsf{meist keine getrennte Prüfsumme für Header und Nutzdaten}, deshalb ist der Header nicht separat überprüfbar
    \end{itemize}
\end{itemize}
\fatsf{Besser wäre: Empfänger meldet Fehler zurück $\rightarrow
$ Negative ACK (NAK)}
\begin{itemize}
    \item Beispiel: \enquote{Ich habe ein fehlerhaftes oder fehlendes Paket erkannt, bitte sende es erneut}
    \item Vorteil:
    \begin{itemize}
        \item Spart Zeit, weil der Sender \fatsf{nicht auf Timeout warten muss}
        \item Besonders bei \fatsf{verzögerungsanfälligen Netzwerken (Internet)} sinnvoll
    \end{itemize}
    \item \fatsf{Aber ist das möglich, wenn der Empfänger das Paket einfach verwirft}
    \begin{itemize}
        \item Ja, \fatsf{indirekt}:
        \begin{itemize}
            \item Wenn das \fatsf{nächste Paket} empfängen wird, erkennt der Empfänger:
            \begin{itemize}
                \item \enquote{Moment, das Paket ist \fatsf{außer der Reihenfolge!}} $\Rightarrow$ Fehler erkannt!
                \item So kann er einen \fatsf{NAK} senden oder auf das verlorene Paket hinweisen
            \end{itemize}
        \end{itemize}
    \end{itemize}
\end{itemize}
\fatsf{Reihenfolgenummern sind unerlässlich}
\begin{itemize}
    \item Um:
    \begin{itemize}
        \item \fatsf{Verlorene Pakete zu erkennen}
        \item Pakete \fatsf{korrekt zuzuordnen}
        \item Daten beim Empfänger wieder \fatsf{in richtige Reihenfolge zu bringen}
        \item In \fatsf{ACKs/NAKs referenzieren} zu können
    \end{itemize}
    \item Besonders kritisch bei \fatsf{IP}, da Pakete \fatsf{nicht zwingend in Reihenfolge} ankommen
\end{itemize}
\fatsf{Letzter Ausweg: Timeout beim Sender}
\begin{itemize}
    \item Wenn keine Bestätigung für ein Paket kommt $\Rightarrow$ \fatsf{Sender wartet Timeout ab} und sendet \fatsf{erneut}
    \item Wichtig:
    \begin{itemize}
        \item Timeout-Dauer ist entscheidend: \fatsf{zu kurz $\rightarrow$ unnötige Wiederholungen, zu lang $\rightarrow$ unnötige Verzögerungen}
    \end{itemize}
\end{itemize}
\subsubsection{Automatic Repeat reQuest (ARQ)}
ARQ = Protokollmechanisumus, bei dem der \fatsf{Sender verlorene oder fehlerhafte Pakete erneut sendet}
\begin{enumerate}
    \item \fatsf{Typ 1: Sotp-And-Wait ARQ}
    \begin{itemize}
        \item Nach jedem Paket wartet der Sender auf ein ACK
        \item Nur \fatsf{ein Paket gleichzeitig in der \enquote{Luft}}
        \item \fatsf{Einfach}, aber \fatsf{sehr ineffizient bei langen Leitungen} (hohe RTT)
        \item Siehe auch: \fatsf{Alternating Bit Protocol (ABP)}
    \end{itemize}
    \item \fatsf{Typ 2: Go-Back-N ARQ}
    \begin{itemize}
        \item Sender kann \fatsf{mehrere Pakete auf einmal} losschicken (Fenstertechnik)
        \item Wenn ein Paket verloren geht oder fehlerhaft ist:
        \begin{itemize}
            \item Empfänger \fatsf{verwirft alle folgenden Pakete}, da sie \enquote{out of order} sind
            \item Sender wird per Timeout oder NACK benachrichtigt und \fatsf{sendet alle ab dem fehlerhaften Paket erneut}
        \end{itemize}
        \item \fatsf{Warum?}
        \begin{itemize}
            \item Früher hatten viele Empfänger nur \fatsf{einen Pufferplatz} $\rightarrow$ konnten keine out-of-order Pakete speichern
        \end{itemize}
    \end{itemize}
    \begin{figure}[h]
        \centering
        \includegraphics[width=0.9\linewidth]{Chapter_4_Images/Go-Back-N ARQ.jpeg}
    \end{figure}
    \clearpage
    \item \fatsf{Typ 3: Selective Reject (Selective Repeat)}
    \begin{itemize}
        \item \fatsf{Effizienter:}
        \begin{itemize}
            \item Empfänger kann mehrere Pakete puffern
            \item Nur das \fatsf{konkrete verlorene Paket wird erneut angefordert}
            \item Sender muss gezielt \fatsf{nur das fehlende Paket erneut senden}
            \item Vorteil:
            \begin{itemize}
                \item \fatsf{Weniger Overhead, bessere Bandbreitenausnutzung}
            \end{itemize}
        \end{itemize}
    \end{itemize}
    \begin{figure}[h]
        \centering
        \includegraphics[width=0.9\linewidth]{Chapter_4_Images/Selective Reject.jpeg}
    \end{figure}
\end{enumerate}
\subsubsection*{Intermezzo: Wie macht es TCP?}
\begin{itemize}
    \item TCP arbeitet \fatsf{nicht mit Paketnummern}, sondern mit \fatsf{Bytenummern}
    \begin{itemize}
        \item Beispiel: Byte 513, Byte 769, ...
    \end{itemize}
    \item Das erlaubt eine \fatsf{feinere Steuerung} und \fatsf{Teilmengenbestätigung}
    \item ABER: Moderne TCP-Implementierungen sind \fatsf{komplex}, da sie \fatsf{Abwärtskompatibilität} wahren müssen
    \begin{itemize}
        \item Deshalb kombinieren sie \fatsf{ACKs mit Bytenummern}, ggf. mit \fatsf{SACK (Selective ACKs)} für selektive Wiederholungen
    \end{itemize}
\end{itemize}
\clearpage
\subsection{Congestion Control}
\subsubsection*{Grundidee: Wo entsteht Überlast (Congestion)?}
Bei jeder Kommunikation im Internet sind \fatsf{mehrere Komponenten beteiligt:}
\begin{enumerate}
    \item \fatsf{Sender (Initiator)}
    \item \fatsf{Netzwerk (Router, Switches, Leitungen etc.)}
    \item \fatsf{Empfänger (Responder)}
\end{enumerate}
Überall dort kann ein \fatsf{Engpass (Bottleneck)} entstehen\\
\subsubsection*{Drei mögliche Bottlenecks:}
\begin{table}[H]
    \centering
    \rowcolors{2}{\thepagecolor}{fgcolor!10!\thepagecolor}
    \renewcommand{\arraystretch}{1.2}
    \begin{tabular}{ll}
        \toprule
        \fatsf{Am Empfänger} & Receiver kann nicht schnell genug lesen oder verarbeiten\\
        \fatsf{Im Kommunikationskanal} & Leitung oder Router sind überlast $\rightarrow$ \fatsf{Netzüberlastung} entsteht\\
        \fatsf{Am Sender (selten)} & Sender limitiert seine Rate selbst, z.B. bei Brandbreitenbegrenzung\\
        \bottomrule 
    \end{tabular}
\end{table}
\begin{figure}[h]
    \centering
    \includegraphics[width=0.9\linewidth]{Chapter_4_Images/Bottleneck.jpeg}
\end{figure}
\clearpage
\subsubsection*{Was ist Congestion (Überlast)?}
\begin{itemize}
    \item \fatsf{Jedes Netzwerk} hat eine \fatsf{begrenzte Kapazität} (Daten pro Sekunde)
    \item Wenn \fatsf{mehr Daten eingespeist} werden, als das Netzwerk verarbeiten kann, entsteht eine \fatsf{Überlastung}
    \item Folge:
    \begin{itemize}
        \item \fatsf{Verzögerungen} (Delays)
        \item \fatsf{Paketverluste} (Drops)
        \item \fatsf{Ineffizienz}, da gleiche Pakete oft mehrmals übertragen werden müssen
        \item \fatsf{Congestive Collapse:} Netzwerk ist aktiv, aber es kommt kaum Daten durch $\rightarrow$ total ineffizient
    \end{itemize}
\end{itemize}
\begin{figure}[h]
    \centering
    \includegraphics[width=0.7\linewidth]{Chapter_4_Images/Congestion.jpeg}
\end{figure}
\clearpage
\subsubsection*{Drei Szenarien zur Veranschaulichung}
\begin{itemize}
    \item \fatsf{Szenario 1 - Router mit unendlichem Puffer, keine Retransmission}
    \begin{itemize}
        \item Zwei Sender und zwei Empfänger kommunizieren über einen Router
        \item Der Router hat \fatsf{unendlich große Puffer}, sodass \fatsf{keine Pakete verloren gehen}
        \item Die \fatsf{eingehende Rate $\lambda_{in}$} kann die \fatsf{Ausgangskapazität C} des Routers überschreiten
        \item Pakete stauen sich im Puffer, was zu \fatsf{wachsenden Verzögerungen} führt
        \item Es entsteht \fatsf{Congestion durch Verzögerung}, nicht durch Verluste
        \item Das Ergebnis: \fatsf{Maximale Dursatzrate C} wird erreicht, aber mit großen Latenzen
    \end{itemize}
    \begin{figure}[h]
        \centering
        \includegraphics[width=0.7\linewidth]{Chapter_4_Images/Szenario 1 P1.jpeg}
    \end{figure}
    \begin{figure}[h]
        \centering
        \includegraphics[width=0.7\linewidth]{Chapter_4_Images/Szenario 1 P2.jpeg}
    \end{figure}
    \clearpage
    \item \fatsf{Szenario 2 - Router mit begrenztem Puffer, mit Retransmission}
    \begin{itemize}
        \item Router hat \fatsf{begrenzten Pufferplatz}; bei Überlast werden Pakete \fatsf{verloren}
        \item Sender erkennen verlorene Pakete (z.B. durch Timeout) und \fatsf{senden diese erneut} (Retransmission)
        \item Die Wiederholungen erhöhen die eingehende Daterate $\lambda_{in}$ zusätzlich
        \item Dadurch entsteht ein \fatsf{Teufelkreis: mehr Last $\rightarrow$ mehr Verluste $\rightarrow$ mehr Retransmissionen $\rightarrow$ noch mehr Last}
        \item Ergebnis: \fatsf{hohe Verzögerungen, verminderte Effizienz}, und Gefahr von \fatsf{Congestion Collapse} trotz physikalischer Kapazitätsgrenze C
    \end{itemize}
    \begin{figure}[h]
        \centering
        \includegraphics[width=0.7\linewidth]{Chapter_4_Images/Szenario 2 P1.jpeg}
    \end{figure}
    \begin{figure}[h]
        \centering
        \includegraphics[width=1\linewidth]{Chapter_4_Images/Szenario 2 P2.jpeg}
    \end{figure}
    \clearpage
    \item \fatsf{Szenario 3: Mehrere Sender, Mehrhops und Retransmissionen}
    \begin{enumerate}
        \item \fatsf{Ausgangslage:}
        \begin{itemize}
            \item Mehrere Sender (z.B. 4 Hosts) senden gleichzeitig
            \item Daten passieren mehrere Router (Multihop)
            \item Router haben begrenzte Pufferkapazität
        \end{itemize}
        \item \fatsf{Problem:}
        \begin{itemize}
            \item Pakete gehen an Routern verloren, wenn Puffer voll sind
            \item Sender erkennen Verluste erst nach Timeout und senden Pakete erneut (Retransmission)
        \end{itemize}
        \item \fatsf{Folge:}
        \begin{itemize}
            \item Retransmissionen erhöhen den Datenverkehr zusätzlich 
            \item Mehr Last führt zu noch mehr Paketverlusten und Verzögerungen
            \item Ressourcen (Bandbreite und Puffer) werden verschwendet
            \item Gefahr von Congestion Collapse steigt
        \end{itemize}
        \item \fatsf{Herausforderung:}
        \begin{itemize}
            \item Staukontrolle ist schwierig, weil Überlast an verschiedenen Punkte entsteht
            \item Kein garantierter Paketreihenfolge oder direkte Verlustmeldung
        \end{itemize}
    \end{enumerate}
    \begin{figure}[h]
        \centering
        \includegraphics[width=0.7\linewidth]{Chapter_4_Images/Szenario 3 P1.jpeg}
    \end{figure}
    \begin{figure}[h]
        \centering
        \includegraphics[width=0.7\linewidth]{Chapter_4_Images/Szenario 3 P2.jpeg}
    \end{figure}
    \clearpage
\end{itemize}
\begin{itemize}
    \item \fatsf{Warum ist Congestion Control nötig?}
    \begin{itemize}
        \item Verhindert Überlastungs-Kettenreaktionen (Snowball-Effekt)
        \item Überlast führt zu Paketverlusten $\rightarrow$ erneute Sendungen $\rightarrow$ noch mehr Last $\rightarrow$ Zusammenbruch
    \end{itemize}
    \item \fatsf{Anpassung der Sendrate:}
    \begin{itemize}
        \item Sender passen ihre Rate an die aktuelle Kapazität des Netzwerks an
        \item Komplex, da viele Router, Weiterleitungsstrategien und Nutzer beteiligt sind
        \item Unterschied zu Flow Control: Flow Control ist lokal (Sender-Empfänger), Congestion Control betrifft das gesamte Netzwerk dazwischen
    \end{itemize}
    \item \fatsf{Ziele der Congestion Control:}
    \begin{itemize}
        \item Hohe Paketlieferung mit kurzen Verzögerungen
        \item Schutz vor kollapsartiger Überlast
        \item Faire Verteilung der Netzwerkressourcen auf alle Datenflüsse
    \end{itemize}
    \item \fatsf{Designoptionen:}
    \begin{itemize}
        \item \fatsf{Open Loop:} Vorab so gestaltet, dass Laufzeitkorrekturen nicht nötig sind
        \item \fatsf{Closed Loop:} Laufzeit-Feedback zur Anpassung der Sendrate
        \begin{itemize}
            \item \fatsf{Explizites Feedback:} Netzwerk signalisiert Überlast direkt
            \item \fatsf{Impliziertes Feedback:} Sender schließt aus Netzwerkverhalten (z.B. fehlende ACKs) auf Überlast
        \end{itemize}
    \end{itemize}
    \item \fatsf{Beispiel:} TCP verwendet implizites Closed-Loop Feedback zur Steuerung der Sendrate
\end{itemize}
\subsubsection*{Mögliche Maßnahmen bei Überlastung}
\begin{itemize}
    \item \fatsf{Optionen zur Überlastvermeidung}
    \begin{enumerate}
        \item \fatsf{Kapazität erhöhen}
        \begin{itemize}
            \item Mehr Links oder Router hinzufügen
            \item \fatsf{Aber:} In der Praxis kurzfristig kaum machbar - teuer und langsam
        \end{itemize}
        \item \fatsf{Reservierung \& Admission Control}
        \begin{itemize}
            \item Neue Datenflüsse nur zulassen, wenn noch genügend Kapazität vorhanden ist
            \item \fatsf{Gilt vor allem für Leitungsvermittlung} (wie bei alten Telefonsystemen)
            \item Funktioniert schlecht im Internet, da Informationen über den aktuellen Zustand selten verfügbar sind
        \end{itemize}
        \item \fatsf{Last reduzieren (feingranular)}
        \begin{itemize}
            \item Quellen verringern ihre Datenrate, \fatsf{ohne} die Verbindung zu beenden
            \item Das funktioniert \fatsf{nur}, wenn sie \fatsf{Feedback} vom Netzwerk bekommen - z.B. durch verlorene Pakete $\rightarrow$ Closed-loop control
        \end{itemize}
    \end{enumerate}
\end{itemize}
\subsubsection*{Klassifikaiton von Maßnahmen}
\begin{itemize}
    \item \fatsf{Wo wird gesteuert?}
    \begin{itemize}
        \item \fatsf{Router-zentriert:} Router trifft Entscheidungen
        \item \fatsf{Host-zentriert:} Endgeräte (Sender/Empfänger) treffen Entscheidungen
        \begin{itemize}
            \item Im Internet fast immer \fatsf{host-zentriert, außer beim Packet-Dropping}
        \end{itemize}
    \end{itemize}
    \item \fatsf{Wie wird die zulässige Datenmenge beschrieben?}
    \begin{itemize}
        \item \fatsf{Window-basiert (wie TCP):}
        \begin{itemize}
            \item Der Sender darf nur ein bestimmtes Fenster an Bytes senden, ohne ACK
        \end{itemize}
        \item \fatsf{Rate-basiert:}
        \begin{itemize}
            \item Eine feste Datenrate (z.B. 100KB/s)
        \end{itemize}
        \item Weitere, seltenere Methoden: z.B. \fatsf{Credit-based}
    \end{itemize}
\end{itemize}
\subsubsection*{Was tun, wenn der Router voll ist?}
\begin{itemize}
    \item \fatsf{Ein neues Paket kommt an, aber der Puffer sit voll}
    \begin{itemize}
        \item Muss \fatsf{ein Paket gedroppt werden} - welches?
        \begin{itemize}
            \item \fatsf{Option 1: Drop Tail}
            \begin{itemize}
                \item Das neue Paket wird einfach verworfen
                \item Vorteil: alte Pakete sind oft \enquote{wertvoller} (z.B. bei Go-Back-N)
                \item \fatsf{Nachteil:} Gefahr der Synchronisation - alle Sender reduzieren gleichzeitig
            \end{itemize}
            \item \fatsf{Option 2: Ein altes Paket verwerfen}
            \begin{itemize}
                \item Besonders sinnvoll bei \fatsf{Multimedia-Streams} $\rightarrow$ neue Daten wichtiger
                \item Oder gezielt ein Paket aus dem gliechen Datenstrom löschen
                \item \fatsf{Aber:} schwer umzusetzen (Overhead)
            \end{itemize}
        \end{itemize}
    \end{itemize}
\end{itemize}
\subsubsection*{Paketverlust = implizites Feedback}
\begin{itemize}
    \item \fatsf{Paketverlust als Signal}
    \begin{itemize}
        \item Wenn ein Paket \fatsf{verloren geht}, kann der Sender (z.B. TCP) das durch fehlendes ACK erkennen
        \item Reaktion: \fatsf{Sende-Rate verringern}
        \item \fatsf{Annahme dabei:} Verluste = Überlast (gilt in drahtgebundenen Netzen oft, aber nicht immer in drahtlosen)
    \end{itemize}
    \item \fatsf{Problem:}
    \begin{itemize}
        \item In drahtlosen Netzwerken: viele Verluste \fatsf{ohne Überlast}, z.B. durch Interferenzen
        \item Lösung: Drahtlose Hops absichern (z.B. durch Link-Layer-Retransmissions)
    \end{itemize}
    \item \fatsf{Open-loop Control:}
    \begin{itemize}
        \item Falls trotz Ressourcenreservierung Pakete an vollen Queues ankommen $\rightarrow$ schlechtes Design
    \end{itemize}
\end{itemize}
\subsubsection*{Proaktive Maßnahmen zur Vermeidung voller Queues}
\begin{itemize}
    \item \fatsf{Idee:}
    \begin{itemize}
        \item Nicht erst reagieren, \fatsf{wenn die Queue voll ist}, sondern schon \fatsf{vorher!}
    \end{itemize}
    \item \fatsf{Router erkennt Überlast frühzeitig, z.B.:}
    \begin{itemize}
        \item Durchschnittliche Queue-Länge über Grenzwert
        \item Auslastung des Ausgangslinks konstant hoch
        \item $\rightarrow$ Router wechselt in den \fatsf{Warnzustand}
    \end{itemize}
\end{itemize}
\begin{figure}[h]
    \centering
    \includegraphics[width=0.4\linewidth]{Chapter_4_Images/Proaktive Maßnahmen zur Vermeidung voller Queues.jpeg}
\end{figure}
\subsubsection*{Proaktive Maßnahme: Choke Packets}
\begin{itemize}
    \item \fatsf{Choke Packet}
    \begin{itemize}
        \item Ein spezielles Paket, das zurück zum Sender geschickt wird: \enquote{Langsamer senden!}
    \end{itemize}
    \item \fatsf{Probleme:}
    \begin{enumerate}
        \item Führt zu \fatsf{noch mehr Datenverkehr}, obwohl das Netz schon überlastet ist
        \item \fatsf{Verzögerung:} Bis der Sender das Feedback erhält, hat er oft schon viele Pakete gesendet (Stichwort: \fatsf{Bandwidth-Delay Product})
    \end{enumerate}
\end{itemize}
\subsubsection*{Proaktive Maßnahme: Warining Bit}
\begin{itemize}
    \item \fatsf{Alternative zu Choke Packets:}
    \begin{itemize}
        \item Router setzt ein \fatsf{Warnbit} im Header der Datenpakete
        \item Der Empfänger kopiert das Bit in die \fatsf{ACKs}
        \item Der Sender sieht das Warnbit in den ACKs $\rightarrow$ verringert seine Sende-Rate
        \item Vorteil: Kein zusätzliches Paket nötig - \fatsf{effizienter}
    \end{itemize}
\end{itemize}
\subsubsection*{Proaktive Maßnahme: Random Early Detection (RED)}
\begin{itemize}
    \item \fatsf{RED-Mechanismus}
    \begin{itemize}
        \item Wenn sich Überlast \fatsf{anbahnt}, werden \fatsf{einige Pakete frühzeitig absichtlich verworfen!}
        \item Klingt brutal, aber:
        \begin{itemize}
            \item Führt dazu, dass Sender \fatsf{frühzeitig drosseln}
            \item So wird verhindert, dass \fatsf{alle Pakete} verloren gehen, wenn es richtig voll wird
        \end{itemize}
    \end{itemize}
    \item \fatsf{Dropprobabilität}
    \begin{itemize}
        \item Beginnt bei 0 (unter \fatsf{MinThreshold})
        \item Steigt \fatsf{linear} zwischen MinThreshold und MaxThreshold
        \item Ab MaxThreshold: \fatsf{100\% Dropprobabilität}
    \end{itemize}
\end{itemize}
\begin{figure}[h]
    \centering
    \includegraphics[width=0.4\linewidth]{Chapter_4_Images/RED.jpeg}
\end{figure}
\subsubsection*{Reaktion auf Feedback}
\fatsf{Was macht der Sender, wenn er Feedback bekommt?}
\begin{itemize}
    \item \fatsf{Rate-basierte Protkolle:}
    \begin{itemize}
        \item Verringern einfach die Sende-Rate (z.B. um Faktor 2)
        \item Frage: Wie und wann wieder erhöhten?
    \end{itemize}
    \item \fatsf{Window-basierte Protkolle (wie TCP):}
    \begin{itemize}
        \item Verringern \fatsf{Congestion Window}
        \item Viele Strategien: z.B. \enquote{Multiplicative Decrease, Additive Increase}
        \item Frage: Wie schnell darf das Fenster wachsen?
        \item Große $\rightarrow$ Burst-Sendungen $\rightarrow$ \fatsf{kann zu neuer Überlast führen}
    \end{itemize}
\end{itemize}
\clearpage
\subsection{Internet UDP}
\subsubsection*{Warum gibt es UDP überhaupt?}
\begin{enumerate}
    \item \fatsf{Kein Verbindungsaufbau = weniger Verzögerung:}
    \begin{itemize}
        \item TCP z.B. muss man erst eine Verbindung aufbauen, was Zeit kostet
        \item UDP sendet einfach drauflos
    \end{itemize}
    \item \fatsf{Einfachheit:}
    \begin{itemize}
        \item Kein Verbindungsstatus wird am Sender oder Empfänger gespeichert
        \item Keine aufwändige Steuerung
    \end{itemize}
    \item \fatsf{Kleiner Header:}
    \begin{itemize}
        \item UDP-Header ist nur 8 Bytes lang (im Gegensatz zu TCP mit 20 Bytes oder mehr)
    \end{itemize}
    \item \fatsf{Keine Staukontrolle (Congestion Control)}
    \begin{itemize}
        \item UDP schickt Pakete so schnell wie möglich
        \item Kein langsamer machen bei Netzwerküberlastung
    \end{itemize}
    \item \fatsf{Warum nicht nur IP?}
    \begin{itemize}
        \item IP liefert Pakete nur bis zum Zielhost
        \item UDP ermöglicht mehrere Anwendungen auf demselben Host (über Ports)
        \item UDP ergänzt IP um eine einfache Multiplexing-Funktion (Ports)
    \end{itemize}
\end{enumerate}
\subsubsection*{Wofür wird UDP oft verwendet?}
\begin{itemize}
    \item \fatsf{Streaming Multimedia:}
    \begin{itemize}
        \item Video, Audio, Live-Übertragungen
        \item Verlust einiger Pakete ist tolerierbar
        \item Wichtig ist, dass Daten schnell ankommen, nicht unbedingt komplett oder in der richtigen Reihenfolge
    \end{itemize}
    \item \fatsf{Andere Beispiele:}
    \begin{itemize}
        \item \fatsf{DNS} (Domain Name System)
        \item \fatsf{SNMP} (Netzwerkmanagement)
        \item Anwendungen, die ihre eigene Fehlerbehandlung machen
    \end{itemize}
\end{itemize}
\subsubsection*{Zuverlässigkeit \& Reihenfolge}
\begin{itemize}
    \item UDP garantiert \fatsf{keine Reihenfolge} und keine Zuverlässigkeit
    \item Wenn Anwendungen zuverlässige oder geordnete Daten wollen, müssen sie selbst \fatsf{Sequneznummern und Fehlerkorrektur} einbauen
    \item Beispiel: TCP ist zuverlässig, UDP nicht - aber man kann mit UDP selbst sowas programmieren, wenn man will
\end{itemize}
\clearpage
\subsubsection*{UDP-Segment Aufbau (Header)}
\begin{figure}[h]
    \centering
    \includegraphics[width=0.5\linewidth]{Chapter_4_Images/UDP Segment.jpeg}
\end{figure}
\subsubsection*{Verbindungslos und Demultiplexing}
\begin{itemize}
    \item UDP benutzt \fatsf{Portnummern}, um Pakete an die richtige Anwendung weiterzuleiten (Demultiplexing)
    \item Beispiel: Ein Rechner kann mehrere UDP-Sockets öffnen, z.B. auf Port 64285, 99111, 99222
    \item Wenn ein UDP-Paket ankommt, schaut das Betriebssystem auf die \fatsf{Zielportnummer} und gibt das Paket an die passende Anwendung (Socket) weiter
    \item Da UDP \fatsf{verbindungslos} ist, gibt es keinen dauerhaften Kanal. Jedes Paket wird einzeln verarbeitet
\end{itemize}
\begin{figure}[h]
    \centering
    \includegraphics[width=0.7\linewidth]{Chapter_4_Images/Connectionless Demultiplexing.jpeg}
\end{figure}
\clearpage
\subsubsection*{Wichtige Punkte}
\begin{itemize}
    \item \fatsf{Socket = IP-Adresse + Port-Nummer} (zweier Tupel, nicht nur Port)
    \item Mehrere Pakete von verschiedenen Quellen mit verschiedenen Quell-IP und Quell-Port können an denselben Ziel-Port (Socket) gehen
    \item Der Quellport dient als \enquote{\fatsf{Rückadresse}}, damit Antworten zum richtigen Sender zurückgehen
\end{itemize}
\subsubsection*{Ziel der UDP-Checksumme}
\begin{itemize}
    \item \fatsf{Fehlererkennung:} Prüfen, ob beim Übertragen eines Segments (UDP-Datenpaket) z.B. \fatsf{Bits gekippt} sind (z.B. $0 \rightarrow 1$ oder $1 \rightarrow$ 0)
    \item Die UDP-Checksumme ist \fatsf{einfach}, aber effektiv gegen viele Übertragungsfehler
\end{itemize}
\subsubsection*{Wie funktioniert die UDP-Checksumme?}
\begin{itemize}
    \item \fatsf{Sender-Seite (vor dem Versand):}
    \begin{enumerate}
        \item \fatsf{Teile das gesamte UDP-Segement} (Header + Daten) in \fatsf{16-Bit-Wörter} auf
        \item \fatsf{Addiere alle 16-Bit-Wörter} miteinander
        \begin{itemize}
            \item Bei Überlauf (also \enquote{Carry out} von links), \fatsf{wrap-around} $\rightarrow$ du addierst ihn \fatsf{wieder dazu}
        \end{itemize}
        \item \fatsf{Nimm das 1-Komplement} (also invertiere jedes Bit: $0 \rightarrow 1$, $1 \rightarrow 0$)
        \item \fatsf{Schreibe dieses Ergebnis in das Checksum-Feld} des UDP-Headers
        \item Versende das Segment
    \end{enumerate} 
    \item \fatsf{Empfänger-Seite (nach Empfang):}
    \begin{enumerate}
        \item Erhalte das UDP-Segment inkl. Checksum-Wert
        \item \fatsf{Addiere alle 16-Bit-Wörter} (inkl. der empfangenen Checksum)
        \item Auch hier: Bei Überlauf - Wrap-around wieder addieren
        \item \fatsf{Ergebnis sollte 0xFFFF sein} (alle Bits = 1)
        \begin{itemize}
            \item Wenn ja $\rightarrow$ vermutlich \fatsf{kein Fehler}
            \item Wenn nein $\rightarrow$ \fatsf{Fehler erkannt}, Segment wird verworfen
        \end{itemize}
    \end{enumerate}
\end{itemize}
\subsubsection*{Beispiel: Internet Checksum}
\begin{figure}[h]
    \centering
    \includegraphics[width=0.5\linewidth]{Chapter_4_Images/Beispiel: Internet Checksuum.jpeg}
\end{figure}
